% =============================================================================
% Chapter 13 -- Causality Experimental Setup
% =============================================================================

\chapter{Causality Experimental Setup}
\label{ch:causality_setup}

This chapter details the experimental methodology for the degradation-and-recovery study: datasets, degradation protocol, recovery protocol, classifiers, and the leakage-safe cross-validation protocol.
The uniform random removal protocol (Chapter~\ref{ch:causality_intro}, Section~\ref{sec:causal:random_removal}) ensures that the degradation sequence introduces no distributional bias, providing the cleanest possible test of the imbalance hypothesis.


% ============================================================================
\section{Datasets}
\label{sec:setup_causal:datasets}

We use ten binary classification datasets drawn from the UCI Machine Learning Repository~\cite{fernandez2018learning}, loaded in a fully balanced configuration ($\IR \approx 1$) by drawing equal-sized random samples from both classes.
Where real dataset files are not available, synthetic analogues are generated using \texttt{sklearn.datasets.make\_classification}~\cite{blagus2013smote} with matching $(n, d)$ parameters; this is a standard practice for controlled ablation studies.

\begin{table}[H]
\centering
\caption{Datasets used in the degradation-recovery study (balanced starting point, $\IR = 1$).}
\label{tab:causal:datasets}
\begin{tabular}{lrrll}
\toprule
\textbf{Dataset} & \textbf{$n$} & \textbf{$d$} & \textbf{Source} & \textbf{Characteristic} \\
\midrule
breast-cancer-wisconsin & 400 &  9 & UCI & Medical, low-dim \\
ionosphere              & 350 & 34 & UCI & Radar, high-dim \\
heart                   & 270 & 13 & UCI & Medical, mid-dim \\
banknote                & 600 &  4 & UCI & Finance, very low-dim \\
sonar                   & 208 & 60 & UCI & Acoustic, very high-dim \\
pima                    & 400 &  8 & UCI & Medical, low-dim \\
haberman                & 200 &  3 & UCI & Survival, minimal features \\
glass                   & 146 &  9 & UCI & Materials, low-dim \\
vehicle                 & 400 & 18 & UCI & Sensor, mid-dim \\
credit                  & 400 & 23 & UCI & Finance, mid-dim \\
\bottomrule
\end{tabular}
\end{table}

The total balanced sample size $n$ ranges from 146 (glass) to 600 (banknote), covering low-, medium-, and high-dimensional settings ($d \in [3, 60]$) representative of real-world imbalanced learning benchmarks~\cite{fernandez2018learning, krawczyk2016learning}.


% ============================================================================
\section{Degradation Protocol}
\label{sec:setup_causal:degradation}

Starting from the balanced dataset, we apply the uniform random removal procedure with the following parameters:

\begin{itemize}[leftmargin=*]
    \item \textbf{Step size:} 1\% of the original minority class removed per step, yielding approximately 80 degradation steps per dataset~\cite{weiss2003effect}.
    \item \textbf{Stop condition:} 80\% removed (20\% of original minority remaining, $\IR \approx 4$).
    This stop point was chosen to match real-world mild-to-moderate imbalance~\cite{he2009learning}.
    \item \textbf{Majority class:} held fixed throughout Phase~1, isolating the minority-removal effect~\cite{japkowicz2002class}.
    \item \textbf{Selection rule:} at each step, a single minority instance is selected uniformly at random for removal.
\end{itemize}

At each step, a single minority instance is selected uniformly at random for removal.
This design ensures no distributional criterion biases the removal sequence, providing the cleanest test of pure imbalance effects.
Because every instance is equally likely to be removed, the degradation protocol avoids confounding the effect of reduced minority count with any systematic bias introduced by a structured removal policy.


% ============================================================================
\section{Recovery Protocol}
\label{sec:setup_causal:recovery}

Phase~2 begins at the most-degraded state (20\% minority remaining, $\IR \approx 4$) and applies each oversampler incrementally:

\begin{itemize}[leftmargin=*]
    \item \textbf{Target sequence:} minority class size increases linearly from $n_{\min}^{(80)}$ to $n_{\min}^{(0)}$ in approximately 80 steps of 1\% each, matching the degradation granularity.
    \item \textbf{Oversampler configuration:} at each recovery step, the oversampler's \texttt{sampling\_strategy} is set to the ratio $n_{\text{target}} / n_{\text{maj}}$, generating only the increment needed for that step rather than a full resample.
    \item \textbf{Oversamplers evaluated:}
    \begin{enumerate}[leftmargin=*]
        \item No-op (lower bound)
        \item Random Oversampling (ROS)~\cite{fernandez2018learning}
        \item SMOTE~\cite{chawla2002smote}
        \item Borderline-SMOTE~\cite{han2005borderline}
        \item ADASYN~\cite{he2008adasyn}
        \item K-Means SMOTE~\cite{douzas2018improving}
        \item GVM-CO (Chapter~\ref{ch:proposed_methods})
    \end{enumerate}
\end{itemize}

The inclusion of GVM-CO directly connects Part~V to Part~IV: the causal study serves as an independent evaluation of GVM-CO's recovery capability in a controlled setting, separate from the benchmark comparisons reported in Chapter~\ref{ch:results}.


% ============================================================================
\section{Classifiers}
\label{sec:setup_causal:classifiers}

Five classifiers spanning the main inductive bias categories are evaluated~\cite{ling2004exploration}: \textbf{RF} (100 trees)~\cite{breiman2001random}, \textbf{DT} (CART, Gini)~\cite{breiman1984cart}, \textbf{KNN} ($k=5$)~\cite{cover1967nearest}, \textbf{LR} (max\_iter$=1000$)~\cite{fernandez2018learning}, and \textbf{MLP} (64-32 hidden units)~\cite{haixiang2017learning}.


% ============================================================================
\section{Leakage-Safe Cross-Validation Protocol}
\label{sec:setup:leakage}

We adopt a strict four-rule protocol to prevent all three forms of data leakage identified by Kaufman et al.~\cite{kaufman2012leakage} and Cawley and Talbot~\cite{cawley2010over}.
Five-fold stratified cross-validation~\cite{kohavi1995study} is used throughout.

\begin{enumerate}[leftmargin=*, itemsep=4pt]
    \item \textbf{Stratified split on pre-oversample data.}
    The $k$-fold stratified split is computed on the \emph{original, non-oversampled} dataset~\cite{cawley2010over}.
    Fold indices are never recomputed after oversampling, preventing the oversampler from influencing which instances appear in which fold.

    \item \textbf{Oversampling inside folds only.}
    For each fold, oversampling is applied exclusively to the training partition~\cite{cawley2010over}.
    The test partition is always original, non-synthetic data.
    No synthetic point ever appears in a test fold.
    This is the single most important rule~\cite{johnson2019survey}; its violation artificially inflates AUC by 2--5\% on typical imbalanced benchmarks~\cite{kaufman2012leakage}.

    \item \textbf{Scaler fitted on training partition only.}
    \texttt{StandardScaler} is fitted on the (already-oversampled) training fold and applied to both training and test folds.
    The test fold's feature scale is never used to fit the scaler, preventing scale leakage~\cite{kaufman2012leakage}.

    \item \textbf{Fresh oversampler per fold.}
    A new oversampler instance is created for each fold, ensuring no cross-fold state (fitted distributions, $k$-NN graphs) is shared~\cite{cawley2010over}.
\end{enumerate}

These four rules are encoded directly in the experiment implementation as \texttt{LeakageSafeCV}.
This protocol prevents a form of evaluation contamination that is common in the imbalanced learning literature: many published results showing large improvements from oversampling are partially attributable to leakage rather than genuine improvements, as demonstrated by Kaufman et al.~\cite{kaufman2012leakage}.
The quantified leakage effect (2--5\% AUC inflation) motivates our strict protocol.


% ============================================================================
\section{Evaluation Protocol}
\label{sec:setup_causal:evaluation}

\begin{itemize}[leftmargin=*]
    \item \textbf{Cross-validation:} 5-fold stratified~\cite{kohavi1995study}, repeated independently at each step in both phases.
    \item \textbf{Reported value:} mean across folds for each metric~\cite{kohavi1995study}.
    \item \textbf{Aggregation across datasets:} macro-average (unweighted mean) across the ten datasets, following standard practice in oversampling benchmarks~\cite{fernandez2018learning}.
    \item \textbf{Metrics:} AUC-ROC, F1-Score, G-Mean, Balanced Accuracy, Sensitivity, Specificity --- the same six metrics used throughout Parts~II--IV.
    \item \textbf{Random seed:} all experiments use \texttt{random\_state = 42} throughout for full reproducibility.
\end{itemize}
